\documentclass[../../../main.tex]{subfiles}
\begin{document}
A generic electromagnetic field is the superposition of many  single frequency monochromatic electromagnetic waves.
Calling $\omega_k$ the angular frequency of the monochromatic wave with wavenumber $\mathbf{k}$, the total energy $H$ of the electromagnetic field is
\begin{equation*}
  H =  \int d^3 \mathbf{r} \left( \frac{\epsilon_0}{2} \mathbf{E}^2 + \frac{1}{2\mu_0} \mathbf{B}^2 \right)                                                                        \\
\end{equation*}
or the semi-classical version
\begin{equation*}
  H = \sum_{\mathbf{k},s} \hbar \omega_k n_{\mathbf{k}s}
\end{equation*}
where $n_{\mathbf{k}s}$ is the number of photons with wavevector $\mathbf{k}$ and polarization $s$.
Taking into account the second quantization of the field
\begin{equation*}
  \hat{H }=\sum_{\mathbf{k}s} \hbar \omega_k \left( \hat{N }_{\mathbf{k} s }+ \frac{1 }{2} \right)
\end{equation*}
with $\hat{N}_{\mathbf{k}s}$ is the quantum number operator
\begin{align*}
  \hat{N}_{\mathbf{k},s} \ket{\dots n_{\mathbf{k}s} \dots}= n_{\mathbf{k}s} \ket{\dots n_{\mathbf{k}s} \dots}
\end{align*}
Each Fourier mode $(\mathbf{k },s)$ acts as a quantum oscillator.
While classical modes correspond to continuous amplitudes, quantum modes have discrete excitations (photons) and zero-point fluctuations.

In classical electrodynamics, the electromagnetic wave can be written in terms of a scalar potential $\phi(\mathbf{r}, t)$ and vector potential $\mathbf{A}(\mathbf{r},t)$
\begin{equation*}
  E = -\nabla \phi - \frac{\partial \mathbf{A}}{\partial t},\qquad
  B = \nabla \land \mathbf{A}
\end{equation*}
This field is of course continuous.
Quantum field treatment however state that the field is quantized with annihilation and creation operator association with it
\begin{align*}
  \hat{\mathbf{E}}(\mathbf{r}, t) & =  i \sum_{\mathbf{k},s} \sqrt{\frac{\hbar \omega_k}{2 \epsilon_0 V}} \left[ \hat{a}_{\mathbf{k}s} e^{i(\mathbf{k} \cdot \mathbf{r} - \omega_k t)} - \hat{a}_{\mathbf{k}s}^{\dagger} e^{-i(\mathbf{k} \cdot \mathbf{r} - \omega_k t)} \right] \boldsymbol{\epsilon}_{\mathbf{k}s}                                        \\
  \hat{\mathbf{B}}(\mathbf{r}, t) & =  \sum_{\mathbf{k},s} \sqrt{\frac{\hbar}{2 \epsilon_0 \omega_k V}} \left[ \hat{a}_{\mathbf{k}s} e^{i(\mathbf{k} \cdot \mathbf{r} - \omega_k t)} - \hat{a}_{\mathbf{k}s}^{\dagger} e^{-i(\mathbf{k} \cdot \mathbf{r} - \omega_k t)} \right] i \frac{\mathbf{k}}{|\mathbf{k}|} \wedge \boldsymbol{\epsilon}_{\mathbf{k}s}
\end{align*}
As with other quantum field theory, photon is the excitation of the quantum electromagnetic field.
Aside, my derivation yield
\begin{equation*}
  \hat{\mathbf{B}}(\mathbf{r}, t) =  \sum_{\mathbf{k},s} \sqrt{\frac{\hbar \omega_k}{2 V} \mu_0} \left[ \hat{a}_{\mathbf{k}s} e^{i(\mathbf{k} \cdot \mathbf{r} - \omega_k t)} - \hat{a}_{\mathbf{k}s}^{\dagger} e^{-i(\mathbf{k} \cdot \mathbf{r} - \omega_k t)} \right] i \mathbf{\hat{k}} \times \boldsymbol{\epsilon}_{\mathbf{k}s}
\end{equation*}
Different constant??

\subsection{The Second Quantization of Light}
The first quantization of light refer to the quantization of electromagnetic energy as photon, while the second quantization of light refer to the quantization of the electromagnetic waves.
This can be achieved in the following steps
\begin{enumerate}
  \item Use Coulomb gauge
  \item Decompose the field
  \item Hamiltonian formulation
  \item Promote the variables to operators
\end{enumerate}

\subsubsection{Coulomb gauge.}
Gauge transformation can be applied to the electromagnetic potential since they do not define electromagnetic wave uniquely
\begin{equation*}
  \phi \rightarrow \phi' = \phi + \frac{\partial }{\partial t} \Lambda(\mathbf{r },t)\qquad
  A \rightarrow A' = A - \nabla \Lambda(\mathbf{r },t)
\end{equation*}
We then use the Coulomb gauge such that
\begin{equation*}
  \nabla  \cdot \mathbf{A } =0
\end{equation*}
From the Coulomb gauge and Gauss law $\nabla \cdot \mathbf{E}=\rho/\epsilon_0$ in vacuum, we also obtain the Laplace equation by applying gradient into it
\begin{equation*}
  \nabla ^2 \phi = 0
\end{equation*}
Imposing that the scalar potential is zero at infinity, this Laplace’s equation has the unique solution
\begin{equation*}
  \phi(\mathbf{r},t)=0
\end{equation*}
and consequently
\begin{equation*}
  \mathbf{E} = - \frac{\partial \mathbf{A}}{\partial t} \qquad \mathbf{B} = \nabla \land \mathbf{A}
\end{equation*}

\subsubsection{Quantum electromagnetic field.}
We can expand the vector potential $\mathbf{A}(r, t)$ as a Fourier series of monochromatic plane waves.
\begin{equation*}
  \mathbf{A}(\mathbf{r}, t) = \sum_{\mathbf{k},s} A_{\mathbf{k} s}(t) \frac{e^{i \mathbf{k} \cdot \mathbf{r}}}{\sqrt{V}} \boldsymbol{\epsilon}_{\mathbf{k} s}
  =  \sum_{\mathbf{k},s} \left[ A_{\mathbf{k}s}(t) \frac{e^{i\mathbf{k} \cdot \mathbf{r}}}{\sqrt{V}} + A_{\mathbf{k}s}^*(t) \frac{e^{-i\mathbf{k} \cdot \mathbf{r}}}{\sqrt{V}} \right] \boldsymbol{\epsilon}_{\mathbf{k}s}
\end{equation*}
Note that the first representation count over all possible $\mathbf{k}$: both $\mathbf{k}$ and $-\mathbf{k}$ are included.
Meanwhile the second only the representation $\left\{ \mathbf{k},-\mathbf{k} \right\} $ are included.


With this, we can write the  quantized electric field
\begin{equation*}
  \mathbf{E}(\mathbf{r}, t) = -\sum_{\mathbf{k},s} \left[ \dot{A}_{\mathbf{k}s}(t)\frac{e^{i \mathbf{\mathbf{k}} \cdot \mathbf{r}}}{\sqrt{V}} + \dot{A}_{\mathbf{k}s}^*(t)\frac{e^{-i \mathbf{\mathbf{k}} \cdot \mathbf{r}}}{\sqrt{V}} \right] \boldsymbol{\epsilon}_{\mathbf{k}s}
\end{equation*}
and magnetic field
\begin{align*}
  \mathbf B(\mathbf r,t)
   & = \nabla \times \mathbf A(\mathbf r,t) \\
   & = \sum_{\mathbf k,s}
  \left[
    A_{\mathbf k s}(t)\frac{\nabla \times \left(e^{i\mathbf k\cdot\mathbf r}\boldsymbol{\epsilon}_{\mathbf k s}\right)}{\sqrt V}
    + A_{\mathbf k s}^*(t)\frac{\nabla \times \left(e^{-i\mathbf k\cdot\mathbf r}\boldsymbol{\epsilon}_{\mathbf k s}\right)}{\sqrt V}
  \right]                                   \\
   & = \sum_{\mathbf k,s}
  \left[
    A_{\mathbf k s}(t)\frac{(\nabla e^{i\mathbf k\cdot\mathbf r})\times\boldsymbol{\epsilon}_{\mathbf k s}}{\sqrt V}
    + A_{\mathbf k s}^*(t)\frac{(\nabla e^{-i\mathbf k\cdot\mathbf r})\times\boldsymbol{\epsilon}_{\mathbf k s}}{\sqrt V}
  \right]                                   \\
   & = \sum_{\mathbf k,s}
  \left[
    A_{\mathbf k s}(t)\frac{(i\mathbf k)e^{i\mathbf k\cdot\mathbf r}\times\boldsymbol{\epsilon}_{\mathbf k s}}{\sqrt V}
    + A_{\mathbf k s}^*(t)\frac{(-i\mathbf k)e^{-i\mathbf k\cdot\mathbf r}\times\boldsymbol{\epsilon}_{\mathbf k s}}{\sqrt V}
  \right]                                   \\
   & = \sum_{\mathbf k,s}
  \left[ A_{\mathbf{k } s}(t)\frac{e^{i \mathbf{k} \cdot \mathbf{r}}}{\sqrt{V}} - A_{\mathbf{k } s}^*(t)\frac{e^{-i \mathbf{k} \cdot \mathbf{r}}}{\sqrt{V}} \right]
  i\mathbf k \times \boldsymbol{\epsilon}_{\mathbf k s}
\end{align*}
where we have used $\nabla \times (f \mathbf{v}) = (\nabla f) \times \mathbf{v}$ for spatial constant $\mathbf{v}$ such that $\nabla  \times \mathbf{v}=0$.

\subsubsection{Quantum electromagnetic field Hamiltonian.}
A familiar result of electromagnetism is that the classical energy of the electromagnetic field in vacuum is given by
\begin{align*}
  H & =  \int d^3 \mathbf{r} \left( \frac{\epsilon_0}{2} \mathbf{E}^2 + \frac{1}{2\mu_0} \mathbf{B}^2 \right)                                                                        \\
  H & =  \int d^3 \mathbf{r} \left( \frac{\epsilon_0}{2} \left( \frac{\partial \mathbf{A}}{\partial t} \right)^2 + \frac{1}{2\mu_0} \left( \nabla \land \mathbf{A} \right)^2 \right)
\end{align*}
To evaluate the Hamiltonian, we evaluate the time derivative and the curl term.
The time derivative is easy enough
\begin{equation*}
  \frac{\partial \mathbf{A}}{\partial t} = \sum_{\mathbf{k},s} \left[ \frac{\dot{A}_{\mathbf{k} s} e^{i \mathbf{k} \cdot \mathbf{r}}}{\sqrt{V}} + \frac{\dot{A}^*_{\mathbf{k} s} e^{-i \mathbf{k} \cdot \mathbf{r}}}{\sqrt{V}} \right] \epsilon_{\mathbf{k} s}
\end{equation*}
To evaluate the curl term use $\nabla \times (f \mathbf{v}) = (\nabla f) \times \mathbf{v}$ again
\begin{align*}
  \nabla \times \mathbf{A} & =  \sum_{\mathbf{k},s} \left[ A_{\mathbf{k}s}(t) \frac{1}{\sqrt{V}} \nabla \times  e^{i\mathbf{k} \cdot \mathbf{r}} + A_{\mathbf{k}s}^*(t) \frac{1}{\sqrt{V}}\nabla \times e^{-i\mathbf{k} \cdot \mathbf{r}} \right] \boldsymbol{\epsilon}_{\mathbf{k}s}                                             \\
  \nabla \times \mathbf{A} & =  \sum_{\mathbf{k},s} \left[ A_{\mathbf{k} s} \frac{i \mathbf{k} \times \boldsymbol{\epsilon}_{\mathbf{k} s}}{\sqrt{V}} e^{i \mathbf{k} \cdot \mathbf{r}} - A^*_{\mathbf{k} s} \frac{i \mathbf{k} \times \boldsymbol{\epsilon}_{\mathbf{k} s}}{\sqrt{V}} e^{-i \mathbf{k} \cdot \mathbf{r}} \right]
\end{align*}

Now the time derivative term reads
\begin{multline*}
  \left( \frac{\partial \mathbf A}{\partial t} \right)^2
  = \sum_{\mathbf k,s} \sum_{\mathbf k', s'} \bigg[
    \dot A_{\mathbf k s} \dot A_{\mathbf k' s'} \frac{e^{i (\mathbf k + \mathbf k') \cdot \mathbf r}}{V}
    + \dot A_{\mathbf k s} \dot A_{\mathbf k' s'}^* \frac{e^{i (\mathbf k - \mathbf k') \cdot \mathbf r}}{V}\\
    + \dot A_{\mathbf k s}^* \dot A_{\mathbf k' s'} \frac{e^{-i (\mathbf k - \mathbf k') \cdot \mathbf r}}{V}
    + \dot A_{\mathbf k s}^* \dot A_{\mathbf k' s'}^* \frac{e^{-i (\mathbf k + \mathbf k') \cdot \mathbf r}}{V}
    \bigg] (\boldsymbol\epsilon_{\mathbf k s} \cdot \boldsymbol\epsilon_{\mathbf k' s'})
\end{multline*}
Recall Fourier transform of 3d Dirac delta
\begin{equation*}
  \int_{\mathbb{R}^3} d^3 \mathbf{r} \, e^{i(\mathbf{k}-\mathbf{k}') \cdot \mathbf{r}} = (2\pi)^3 \delta^3(\mathbf{k}-\mathbf{k}')
\end{equation*}
with continuous $\mathbf{k}$.
Discrete $\mathbf{k}=2 \pi/L(n_x, n_y, n_z), n_i \in \mathbb{Z}$ takes place in system is enclosed in a finite box of volume $V=L^3$ turns the integral into
\begin{equation*}
  \int_V d^3 \mathbf{r} \;e^{i(\mathbf{k}-\mathbf{k}') \cdot r} = V \delta_{\mathbf{k},\mathbf{k}'}
\end{equation*}
Suppose earlier we pick positive value of $\mathbf{k}$.
The term of exponent containing $\mathbf{k} + \mathbf{k}$ vanishes because by the Dirac delta
\begin{equation*}
  \int_V d^3 \mathbf{r} \;e^{i(\mathbf{k}+\mathbf{k}') \cdot r} = V \delta_{\mathbf{k},-\mathbf{k}'}
\end{equation*}
only the terms $\mathbf{k}=-\mathbf{k}'$, which are outside our summation range, survive.
As such, the integral evaluate into
\begin{equation*}
  \int d^3 \mathbf{r} \left( \frac{\partial \mathbf{A}}{\partial t} \right)^2  = \sum_{\mathbf{k},s}  \left( \dot{A}_{\mathbf{k} s}^{*} \dot{A}_{\mathbf{k} s} + \dot{A}_{\mathbf{k} s} \dot{A}_{\mathbf{k} s}^{*} \right)
\end{equation*}
Also, the polarization orthogonality gives $\boldsymbol{\epsilon}_{\mathbf{k} s} \cdot \boldsymbol{\epsilon}_{\mathbf{k} s'} =\delta_{ss'}$.
Vector potential obeys the wave equation whose solution is $A_{\mathbf{k}s}(t)=A_{\mathbf{k}s}(0)e^{i\omega_kt}$.
Thus $\dot{A}_{\mathbf{k}s}=-i\omega_kA_{\mathbf{k}s}$ and
\begin{equation*}
  \frac{\epsilon_0 }{2 }\int d^3 \mathbf{r} \left( \frac{\partial \mathbf{A}}{\partial t} \right)^2
  =\frac{1 }{2} \sum_{\mathbf{k},s}  \omega_k^2 \left( {A}_{\mathbf{k} s}^{*} A_{\mathbf{k} s} + A_{\mathbf{k} s} A_{\mathbf{k} s}^{*} \right)
\end{equation*}

Next the curl term
\begin{multline*}
  (\nabla \times \mathbf A)^2
  = \sum_{\mathbf k,s} \sum_{\mathbf k',s'} \bigg[
    A_{\mathbf k s} A_{\mathbf k' s'} \frac{-(\mathbf k \times \boldsymbol\epsilon_{\mathbf k s}) \cdot (\mathbf k' \times \boldsymbol\epsilon_{\mathbf k' s'})}{V} e^{i (\mathbf k + \mathbf k') \cdot \mathbf r}\\
    + A_{\mathbf k s} A_{\mathbf k' s'}^* \frac{(\mathbf k \times \boldsymbol\epsilon_{\mathbf k s}) \cdot (\mathbf k' \times \boldsymbol\epsilon_{\mathbf k' s'})}{V} e^{i (\mathbf k - \mathbf k') \cdot \mathbf r} \\
    + A_{\mathbf k s}^* A_{\mathbf k' s'} \frac{(\mathbf k \times \boldsymbol\epsilon_{\mathbf k s}) \cdot (\mathbf k' \times \boldsymbol\epsilon_{\mathbf k' s'})}{V} e^{-i (\mathbf k - \mathbf k') \cdot \mathbf r}\\
    + A_{\mathbf k s}^* A_{\mathbf k' s'}^* \frac{-(\mathbf k \times \boldsymbol\epsilon_{\mathbf k s}) \cdot (\mathbf k' \times \boldsymbol\epsilon_{\mathbf k' s'})}{V} e^{-i (\mathbf k + \mathbf k') \cdot \mathbf r}
    \bigg].
\end{multline*}
Pretty long, but that's fine.
The cross term vanishes again I hope\dots
\begin{align*}
  \int d^3 \mathbf{r}  \; (\nabla \times \mathbf A)^2
   & =  \sum_{\mathbf{k},s}  \left( \dot{A}_{\mathbf{k} s}^{*} \dot{A}_{\mathbf{k} s} + \dot{A}_{\mathbf{k} s} \dot{A}_{\mathbf{k} s}^{*} \right) |\mathbf{k }\times \boldsymbol{\epsilon}_{\mathbf{k}s} |^2 \\
   & =  \sum_{\mathbf{k},s} k^2 \left( \dot{A}_{\mathbf{k} s}^{*} \dot{A}_{\mathbf{k} s} + \dot{A}_{\mathbf{k} s} \dot{A}_{\mathbf{k} s}^{*} \right)
\end{align*}
Yup, under Dirac delta it does.
Do note that $|\mathbf{f} \times \mathbf{v}|=|\mathbf{f}| |\mathbf{v}| \sin \theta$ and that the polarization vector is perpendicular to the wave vector.
Using $\omega_k = ck$ and $c^2=1/\epsilon_0\mu_0$ to write $\epsilon_0\omega_k^2=k^2/\mu_0$, the integral turns
\begin{equation*}
  \frac{1 }{2\mu_0} \int d^3 \mathbf{r}  \; (\nabla \times \mathbf A)^2
  =\frac{1 }{2} \sum_{\mathbf{k},s}  \omega_k^2 \left( {A}_{\mathbf{k} s}^{*} A_{\mathbf{k} s} + A_{\mathbf{k} s} A_{\mathbf{k} s}^{*} \right)
\end{equation*}

Therefore, the  Hamiltonian for both field
\begin{equation*}
  H  = \sum_{\mathbf{k},s}  \omega_k^2 \left( {A}_{\mathbf{k} s}^{*} A_{\mathbf{k} s} + A_{\mathbf{k} s} A_{\mathbf{k} s}^{*} \right)
\end{equation*}
In terms of dimensionless complex coefficients
\begin{equation*}
  A_{\mathbf{k}s}(t)=\sqrt{\frac{\hbar }{2\epsilon_0 \omega_k }} a_{\mathbf{k }s}(t)\quad
  A ^*_{\mathbf{k}s}(t)=\sqrt{\frac{\hbar }{2\epsilon_0 \omega_k }} a_{\mathbf{k }s} ^*(t)
\end{equation*}
we can write
\begin{equation*}
  H = \sum_{\mathbf{k},s} \frac{\hbar \omega_{k}}{2} \left( a_{\mathbf{k} s}^* a_{\mathbf{k} s} + a_{\mathbf{k} s} a_{\mathbf{k} s}^* \right)
\end{equation*}
Instead of using the complex amplitudes, one can introduce the real variables
\begin{equation*}
  q_{ks}(t) = \sqrt{\frac{2\hbar}{\omega_k}}\frac{1}{2} \left[ a_{ks}(t) + a_{ks}^*(t) \right]  \quad
  p_{ks}(t) = \frac{\sqrt{2\hbar \omega_k}  }{2i} \left[ a_{ks}(t) - a_{ks}^*(t) \right]
\end{equation*}
such that
\begin{equation*}
  H = \sum_{\mathbf{ k} s}  \left( \frac{p_{\mathbf{k}s}^{2}}{2} + \frac{1}{2} \omega_{k}^{2} q_{\mathbf{k} s}^{2} \right)
\end{equation*}

\subsubsection{Operator promotion.}
The next step is to promote the complex coefficients into ladder operator satifying $[\hat{a}_{\mathbf{k} s}, \hat{a}_{\mathbf{k}'s'}^\dagger] = \delta_{\mathbf{k},\mathbf{k}'} \delta_{s,s'}$
\begin{equation*}
  a_{\mathbf{k} s} \rightarrow \hat{a}_{\mathbf{k}s} \qquad
  a^*_{\mathbf{k}s} \rightarrow \hat{a}^\dagger_ {\mathbf{k}s}
\end{equation*}
The quantum Hamiltonian is thus given by
\begin{equation*}
  \hat{H} = \sum_{\mathbf{k},s} \frac{\hbar \omega_{k}}{2} \left( \hat{a}_{\mathbf{k} s}^\dagger \hat{a}_{\mathbf{k} s} + \hat{a}_{\mathbf{k} s} \hat{a}_{\mathbf{k} s}^\dagger \right)
\end{equation*}
Using the commutator relation $\hat{a}_{\mathbf{k} s}\hat{a}_{\mathbf{k}s}^\dagger= \hat{a}_{\mathbf{k}s}^\dagger\hat{a}_{\mathbf{k} s} +1 $
\begin{equation*}
  \hat{H} = \sum_{\mathbf{k},s} \frac{\hbar \omega_{k}}{2} \left( 2\hat{a}_{\mathbf{k} s}^\dagger \hat{a}_{\mathbf{k} s} + 1\right)
  = \sum_{\mathbf{k},s} \hbar \omega_{k} \left( \hat{a}_{\mathbf{k} s}^\dagger \hat{a}_{\mathbf{k} s} + \frac{1 }{2}\right)
\end{equation*}
Or, using number operator which counts the number of photons in the single-particle in state $\ket{\mathbf{k }s}$
\begin{equation*}
  \hat{H}= \sum_{\mathbf{k},s} \hbar \omega_{k} \left( \hat{N}_{\mathbf{k} s} + \frac{1 }{2}\right)
\end{equation*}

In 1927 Paul Dirac performed the quantization of the classical Hamiltonian by promoting the real coordinates $q_{\mathbf{k }s}$ and the real momenta $p_{\mathbf{k }s}$ to operators satifying $[\hat{q}_{\mathbf{k} s}, \hat{p}_{\mathbf{k}'s'}^\dagger] =i \hbar  \delta_{\mathbf{k},\mathbf{k}'} \delta_{s,s'}$
\begin{equation*}
  q_{\mathbf{k} s} \rightarrow \hat{q}_{\mathbf{k}s} \qquad
  p^*_{\mathbf{k}s} \rightarrow \hat{p}^\dagger_ {\mathbf{k}s}
\end{equation*}
The Hamiltonian is given by
\begin{equation*}
  \hat{H} = \sum_{\mathbf{ k}, s}  \left( \frac{\hat{p}_{\mathbf{k}s}^{2}}{2} + \frac{1}{2} \omega_{k}^{2} \hat{q}_{\mathbf{k} s}^{2} \right)
\end{equation*}
Annihilation and creation operators can be expressed in terms of these operator
\begin{equation*}
  \hat{a}_{\mathbf{k } s} = \sqrt{\frac{\omega_k}{2\hbar}} \left( \hat{q}_{\mathbf{k } s} + \frac{i}{\omega_k} \hat{p}_{\mathbf{k } s} \right) \quad
  \hat{a}_{\mathbf{k } s} ^\dagger= \sqrt{\frac{\omega_k}{2\hbar}} \left( \hat{q}_{\mathbf{k } s} - \frac{i}{\omega_k} \hat{p}_{\mathbf{k } s} \right)
\end{equation*}

Next we move to the electric field.
Since $\dot{A}_{\mathbf{k}s}=-i\omega_kA_{\mathbf{k}s}$, the electric field now reads
\begin{equation*}
  \mathbf{E}(\mathbf{r}, t) = i\sum_{\mathbf{k},s}\omega_k \left[ {A}_{\mathbf{k}s}(t)\frac{e^{i \mathbf{\mathbf{k}} \cdot \mathbf{r}}}{\sqrt{V}} - {A}_{\mathbf{k}s}^*(t)\frac{e^{-i \mathbf{\mathbf{k}} \cdot \mathbf{r}}}{\sqrt{V}} \right] \boldsymbol{\epsilon}_{\mathbf{k}s}
\end{equation*}
Introduce the dimensionless coefficients
\begin{equation*}
  \mathbf{E}(\mathbf{r}, t) = i \sum_{\mathbf{k},s} \sqrt{\frac{\hbar \omega_k}{2 \epsilon_0 V}} \left[ a_{\mathbf{k } s}(t) e^{i\mathbf{k} \cdot \mathbf{r}} - a_{ks}^*(t) e^{-i\mathbf{k} \cdot \mathbf{r}} \right] \epsilon_{\mathbf{k } s}.
\end{equation*}
Finally promote the coefficients into operators
\begin{equation*}
  \hat{\mathbf{E}}(\mathbf{r}, t) =  i \sum_{\mathbf{k},s} \sqrt{\frac{\hbar \omega_k}{2 \epsilon_0 V}} \left[ \hat{a}_{\mathbf{k}s} e^{i(\mathbf{k} \cdot \mathbf{r} - \omega_k t)} - \hat{a}_{\mathbf{k}s}^{\dagger} e^{-i(\mathbf{k} \cdot \mathbf{r} - \omega_k t)} \right] \boldsymbol{\epsilon}_{\mathbf{k}s}
\end{equation*}
where we expand the ladder operator definition using $\hat{a}_{\mathbf{k }s }(t)=\hat{a}_{\mathbf{k }s}e^{-\omega_kt}$.

For the magnetic field case, we begin by inserting the dimensionless amplitude
\begin{equation*}
  \mathbf{B}(\mathbf{r},t) = \sum_{\mathbf k,s}   \sqrt{\frac{\hbar}{2 \epsilon_0 \omega_k V}}
  \left[ a_{\mathbf{k } s}(t)e^{i \mathbf{k} \cdot \mathbf{r}} - a_{\mathbf{k } s}^*(t)e^{-i \mathbf{k} \cdot \mathbf{r}} \right]   i\mathbf k \times \boldsymbol{\epsilon}_{\mathbf k s}
\end{equation*}
then simply promote them into operator
\begin{equation*}
  \hat{\mathbf{B}}(\mathbf{r}, t) =  \sum_{\mathbf{k},s} \sqrt{\frac{\hbar}{2 \epsilon_0 \omega_k V}} \left[ \hat{a}_{\mathbf{k}s} e^{i(\mathbf{k} \cdot \mathbf{r} - \omega_k t)} - \hat{a}_{\mathbf{k}s}^{\dagger} e^{-i(\mathbf{k} \cdot \mathbf{r} - \omega_k t)} \right] i \mathbf{k} \times \boldsymbol{\epsilon}_{\mathbf{k}s}
\end{equation*}
Still missing some factor of $1/|\mathbf{k}|$ to make $\mathbf{k}\to \hat{\mathbf{k}}$ in the last term.
Forcing it regardless, yields
\begin{equation*}
  \hat{\mathbf{B}}(\mathbf{r}, t) =  \sum_{\mathbf{k},s} \sqrt{\frac{\hbar \omega_k}{2 V} \mu_0} \left[ \hat{a}_{\mathbf{k}s} e^{i(\mathbf{k} \cdot \mathbf{r} - \omega_k t)} - \hat{a}_{\mathbf{k}s}^{\dagger} e^{-i(\mathbf{k} \cdot \mathbf{r} - \omega_k t)} \right] i \mathbf{\hat{k}} \times \boldsymbol{\epsilon}_{\mathbf{k}s}
\end{equation*}
a different constant.
Convention difference perhaps?

The operators may be written as a sum over all possible $\mathbf{k}$ rather than only independent $\mathbf{k}$, but this is a matter of convention and not a consequence of quantization itself.
Upon quantization, the independent classical amplitudes are promoted to operators.
The classical reality condition  $A_{-\mathbf{k}} = A_{\mathbf{k}}^*$ is replaced by the requirement that the field operator be Hermitian $\hat{\mathbf A}^\dagger = \hat{\mathbf A}$
This Hermiticity is guaranteed by the structure of the operator expansion, where each annihilation operator is paired with its adjoint multiplying the opposite plane wave.
Whether the summation is taken over all $\mathbf{k}$ or only over independent $\mathbf{k}$ is a bookkeeping choice. In either case, the Hermitian structure of the field operator ensures consistency, and the physical content is unchanged.

\subsection{Zero-Point Energy and the Casimir Effect}
In the case of the quantum electromagnetic field, there is an infinite number of harmonic oscillators and the total zero-point energy,
\begin{equation*}
  \braket{0 | H |0 }\equiv E_{\text{vac}}= \sum_{\mathbf{k },s } \frac{1 }{2 }\hbar \omega_k
\end{equation*}
is clearly infinite.
The infinite constant $E_\text{vac}$ is usually eliminated by simply shifting to zero the energy associated to the vacuum state $\ket{0}$.

Introducing the plates alters the allowed modes of the field, producing a vacuum energy $E'_\text{vac}(a)$ that depends on the plate separation $a$. The vacuum energy in free space, $E_\text{vac}$, serves as a reference.
To quantify the effect, one defines the Casimir energy per unit area as
\begin{equation*}
  \mathcal{E}(a) = \frac{E'_\text{vac}(a) - E_\text{vac}}{L^2}= -\frac{\pi^2 }{720 }\frac{\hbar c }{a^3}
\end{equation*}
Experimentally, two parallel conducting plates are observed to attract each other when brought close together.
This behavior reflects  systems evolution toward configurations of minimum energy.
As the plates approach each other, the restriction on electromagnetic modes increases, reducing $E'_\text{vac}(a)$ and making $\mathcal{E}(a)$ more negative.
This decrease in energy corresponds to a more favorable configuration, consistent with the observed attraction.
The Casimir force per unit area is obtained by differentiating this energy with respect to the plate separation,
\begin{equation*}
  \mathcal{F} = - \frac{\partial \mathcal{E}}{\partial a} = - \frac{\pi^2 }{240 }\frac{\hbar c }{a^4}
\end{equation*}
Defining the energy difference as $E'_\text{vac} - E_\text{vac}$ ensures that the sign of the Casimir energy directly reflects the tendency of the system toward lower energy, so that a decrease in separation corresponds naturally to both a decrease in energy and an attractive force.

\subsubsection{Derivation: Casimir energy per unit area.}
Consider a cubic box of volume \( V = L^3 \) with periodic boundary conditions. The allowed wavevectors are discretized as
\begin{equation*}
  k_x = \frac{2\pi n_x}{L}, \quad k_y = \frac{2\pi n_y}{L}, \quad k_z = \frac{2\pi n_z}{L}, \quad n_x, n_y, n_z \in \mathbb{Z}.
\end{equation*}
Each allowed \( k \) occupies a ``cell'' in \( k \)-space of volume
\begin{equation*}
  \Delta k_x \Delta k_y \Delta k_z = \left( \frac{2\pi}{L} \right)^3.
\end{equation*}
In the continoum limit, the sum over discrete modes can be written as
\begin{equation*}
  \sum_{\mathbf{k},s} f(\mathbf{k},s) \rightarrow  \frac{V}{(2\pi)^{3}} 2\int d^{3}\mathbf{k} \; f(\mathbf{k})
\end{equation*}
The zero-point energy of the electromagnetic field in a region of volume $V$ can be written as
\begin{equation*}
  E_\text{vac}=
  V \int \frac{d^{3}\mathbf{k}}{(2\pi)^{3}} \hbar c \sqrt{k_{x}^{2} + k_{y}^{2} + k_{z}^{2}}
\end{equation*}
In particular, considering a region with the shape of a parallelepiped of length $L$ along both $x$ and $y$ and length $a$ along $z$
\begin{align*}
  E_\text{vac} & = \hbar c \int  \frac{L dk_{x}}{2\pi} \int \frac{L dk_{y}}{2\pi} \int  \frac{a dk_{z}}{2\pi} \sqrt{k_{x}^{2} + k_{y}^{2} + k_{z}^{2}}                                                                           \\
               & = \frac{\hbar  c }{(2\pi)^3 }L^2 a \left( 2\pi \int_{0 }^{\infty }dk_{\parallel}\;k_{\parallel}\int dk_{\parallel}\sqrt{k_{\parallel}^2+k_z^2} \right)                                                          \\
               & = \frac{\hbar  c }{(2\pi)^3 }L^2 a \left[ 2\pi \int_{0 }^{\infty }dk_{\parallel}\;k_{\parallel}\frac{\hbar c}{2\pi}\int_0^\infty dk_{\parallel}\sqrt{k_{\parallel}^2+\left( \frac{n \pi}{a} \right) ^2} \right] \\
  E_\text{vac}
               & =  \frac{\hbar c}{2\pi} L^{2} \int_{0}^{\infty} dk_{\parallel}k_{\parallel}\left[\int_{0}^{\infty} dn \sqrt{k_{\parallel}^{2} + \frac{n^{2} \pi^{2}}{a^{2}}}\right]
\end{align*}
where the expression is obtained setting $k_{\parallel} = \sqrt{k_{x}^{2} + k_{y}^{2}}$ and $n = (a/\pi)k_{z}$.

Let us now consider different  case of metallic plates with length $L \gg a$ having parallel faces lying in the $(x,y)$ plane at distance $a$.
The boundary condition of Maxwell equation for tangential components electric field reads
\begin{equation*}
  \mathbf{E}_1^{\parallel} - \mathbf{E}_2^{\parallel}=0
\end{equation*}
Since we use conductor, the electric inside the conductor must be zero, thus the tangential components must be zero $\mathbf{E}^{\parallel}=0$.
The separation also ensures that $k_x$ and $k_y$ remains continuous but the $k_{z}$ is quantized via
\begin{equation*}
  k_{z} = n \pi/a
\end{equation*}
where now $n = 0, 1, 2, \ldots$ is an integer number
This is equivalent to  Dirichlet boundary conditions where the electric field vanishes at the boundaries.
In this case the zero-point energy in the volume $V = L^{2} a$ between the two plates reads
\begin{equation*}
  E'_\text{vac} =
  \frac{\hbar c}{2\pi} L^{2} \int_{0}^{\infty} dk_{\parallel}k_{\parallel} \left( \frac{k_{\parallel}}{2} + \sum_{n=1}^{\infty} \sqrt{k_{\parallel} ^{2}+ \frac{n^{2} \pi^{2}}{a^{2}}}  \right)
\end{equation*}
where for \( k_z = 0 \) the state \(|(k_z, k_y, 0)1\rangle\)  normal to the \( xy \) plane satisfy $\mathbf{E}^{\parallel}=0$, while the state \(|(k_z, k_y, 0)2\rangle\) orthogonal to the \( z \) axis and lies in the $xy$ plane does not, thus does not contribute to the summation.

This mode restriction result in the lowering zero-point energy.
The difference between the unrestricted vacuum energy and the reduced cavity energy generates a force pulling the plates together.
This is the Casimir force.

The net energy per unit surface area
\begin{align*}
  \mathcal{E} & =  \frac{E'_{\text{vac}} - E_{\text{vac}}}{L^2}                                                                                                                                                                                                                          \\
              & =  \frac{\hbar c}{2\pi}  \int_{0}^{\infty} dk_{\parallel}k_{\parallel} \left( \frac{k_{\parallel}}{2} + \sum_{n=1}^{\infty} \sqrt{k_{\parallel} ^{2}+ \frac{n^{2} \pi^{2}}{a^{2}}}  -\int_{0}^{\infty} dn \sqrt{k_{\parallel}^{2} + \frac{n^{2} \pi^{2}}{a^{2}}} \right) \\
              & =  \frac{\hbar c}{2\pi} \left( \frac{\pi}{a } \right) ^{3} \int_{0}^{\infty} d\zeta \zeta \left( \frac{\zeta}{2} + \sum_{n=1}^{\infty} \sqrt{\zeta ^{2}+ n^2}  -\int_{0}^{\infty} dn \sqrt{\zeta^{2} + n^2} \right)                                                      \\
              & =\pi^2 \frac{\hbar c }{2a^3} \left[ \frac{1}{2} A(0) + \sum_{n=1}^{\infty} A(n) - \int_{0}^{\infty} dn\; A(n) \right]
\end{align*}
with \( \zeta = a k_{\parallel}/\pi \) and
\begin{equation*}
  A(n) = \int_{0}^{+\infty} d\zeta \, \zeta \sqrt{\zeta^2 + n^2}
  = \frac{1}{3} \left[ (n^2 + \infty)^{3/2} - n^2 \right]
\end{equation*}
Although divergent, it might be evaluated as
\begin{equation*}
  A(n)=  \frac{1}{2} \int_{n^{2}}^{\infty} \sqrt{u} \; du= \frac{1 }{2 }\frac{2 }{3 }u^{3/2}\bigg|_{n^2 }^{\infty }=
  \frac{1 }{3 }\left( \infty-n^3 \right)
\end{equation*}
with the substitution of $u = \zeta^{2} + n^{2} $ and $\zeta\;d\zeta=du/2$.
The divergent part, however, cancel each other according to Euler-MacLaurin formula for the difference of infinite series and integrals
\begin{multline*}
  \frac{1}{2} A(0) + \sum_{n=1}^{\infty} A(n) - \int_{0}^{\infty} dn A(n) = -\frac{1}{6 \cdot 2!} \frac{dA}{dn}\bigg|_{n=0} \\
  + \frac{1}{30 \cdot 4!} \frac{d^3 A}{dn^3}\bigg|_{n=0} - \frac{1}{42\cdot 6!}  \frac{d^5 A}{dn^5}\bigg|_{n=0} + \cdots
\end{multline*}
Only the third order derivative survive
\begin{equation*}
  \mathcal{E}
  =\pi^2 \frac{\hbar c }{2a^3} \frac{1}{30 \cdot 4!} \frac{d^3 A}{dn^3}\bigg|_{n=0}
  =-\pi^2 \frac{\hbar c }{2a^3} \frac{1}{30 \cdot 4!} 2
  = -\frac{\pi^2 }{720 }\frac{\hbar c }{a^3}
\end{equation*}

\subsubsection{Derivation: Casimir force.}
The attractive force per unit area between the two plates is simply the separation derivative of the energy per unit area
\begin{equation*}
  \mathcal{F}=-\frac{d \mathbb{E}}{da} = \frac{\pi^{2}}{720} \hbar c \cdot \frac{d}{da} \left( a^{-3} \right) = \frac{\pi^{2}}{720} \hbar c \cdot \left( -3 a^{-4} \right) = -\frac{\pi^{2}}{240} \hbar c \cdot a^{-4}
\end{equation*}

\subsection{Linear and Angular Momentum of the Radiation Field}
In addition to the total energy there are other interesting conserved quantities which characterize the classical radiation field. 
They are the total linear momentum
\begin{equation*}
P = \int d^3 \mathbf{r} \; \epsilon_0 \mathbf{E}(\mathbf{r}, t) \wedge \mathbf{B}(\mathbf{r}, t)
\end{equation*}
or 
\begin{equation*}
\hat{P} = \sum_{\mathbf{k },s} \hbar \mathbf{k} \left( \hat{N}_{\mathbf{k} s} + \frac{1}{2} \right) 
= \sum_{\mathbf{k },s} \hbar \mathbf{k} \hat{N}_{\mathbf{k} s}  
\end{equation*}
Also the angular momentum
\begin{equation*}
  \mathbf{J} = \int d^3 \mathbf{r} \; \epsilon_0 \mathbf{r} \wedge \left[ \mathbf{E}(\mathbf{r}, t) \wedge \mathbf{B}(\mathbf{r}, t) \right] 
\end{equation*}
After evaluating the expression, one obtains 
\begin{equation*}
    \mathbf{\hat{J}} = 
    \int d^3 \mathbf{k} \sum_s \hat{a}^\dagger_{\mathbf{k }s}  \left[i \hbar \mathbf{k} \times \nabla_\mathbf{k} \right] \hat{a}^\dagger_{\mathbf{k }s} 
    + \hbar \int d^3 \mathbf{k} \sum_{s,s'} \hat{a}^\dagger_{\mathbf{k }s}  \mathbf{S}_{ss'} \hat{a}_{\mathbf{k }s} 
\end{equation*}
with $\mathbf{S}_{ss'}$is the matrix representation of the spin-1 rotation generators restricted to the transverse polarization subspace.
It contains two contributions: the orbital part
\begin{equation*}
    \hat{\mathbf{J}}_{\text{orb}} = 
        \int d^3 \mathbf{k} \sum_s \hat{a}^\dagger_{\mathbf{k }s}  \left[ i \hbar \mathbf{k} \times \nabla_\mathbf{k} \right] \hat{a}^\dagger_{\mathbf{k }s} 
\end{equation*}
the spin (helicity) part
\begin{equation*}
    J_{\text{spin}} = 
     \hbar \int d^3 \mathbf{k} \sum_{s,s'} \hat{a}^\dagger_{\mathbf{k }s}  \mathbf{S}_{ss'} \hat{a}_{\mathbf{k }s} 
\end{equation*}

The linearly polarized states \(|s\rangle\)  are eigenstates of \( \hat{H} \) and \( \hat{P} \), but they are not  \( \hat{J} \), \( \hat{J}^2 \), and nor of \( \hat{J}_z \). 
\begin{align*}
  \hat{H} \ket{\mathbf{k }s} &=  \hbar \omega_k \ket{\mathbf{k }s}\\
  \mathbf{\hat{P}} \ket{\mathbf{k}s} &= \hbar \mathbf{k }\ket{\mathbf{k}s}
\end{align*}
The eigenstates of both \(\ket{kjm_js}\) do not have a fixed wavevector \( \mathbf{k} \), only  fixed wavenumber \( k = |k| \). 
These states are associated to vector spherical harmonics \( Y_{j,m}(\theta, \phi) \).
\(\ket{kjm_js}\) are not eigenstates of \( \hat{P} \) but they are eigenstates of \( \hat{H} \), \( \hat{J}^2 \), and \( \hat{J}_z \)
\begin{align*}
\hat{H}\ket{kjm_js} &=  c k \ket{kjm_js} \\
\hat{J}^2 \ket{kjm_js}  &=  j(j + 1)\hbar^2 \ket{kjm_js}\\ 
\hat{J}_z \ket{kjm_js} &=  m_j \hbar \ket{kjm_js} \\
\end{align*}
with $m_j=-j,\dots,j$ and $j=1,\dots$.

\subsubsection{Derivation: momentum operator}
The electromagnetic field term reads 
\begin{multline*}
  \mathbf{E }\times \mathbf{B}
  = - \sum_{\mathbf{k },s} \sum_{\mathbf{k }',s'}  \frac{\hbar  }{2V} \sqrt{\frac{\mu_0 }{\epsilon_0 }\omega_k \omega_{k'}}
    \bigg[ \hat{a}_{\mathbf{k }s} \hat{a}_{\mathbf{k }'s'}  e ^{i(\phi_{\mathbf{k }} +\phi_{\mathbf{k }'})} 
  -\hat{a}_{\mathbf{k }s} \hat{a}^\dagger_{\mathbf{k }'s'} e ^{i(\phi_{\mathbf{k }}-\phi_{\mathbf{k }'})} \\ 
   -\hat{a}_{\mathbf{k }s} ^\dagger\hat{a}_{\mathbf{k }'s'} e ^{-i(\phi_{\mathbf{k }}-\phi_{\mathbf{k }'})}  
   +\hat{a}^\dagger_{\mathbf{k }s} \hat{a}^\dagger_{\mathbf{k }'s'} e ^{-i(\phi_{\mathbf{k }}+\phi_{\mathbf{k }'})} \bigg]
  \boldsymbol{\epsilon}_{\mathbf{k}s} \times \left[ \hat{\mathbf{k }'}\times \boldsymbol{\epsilon}_{\mathbf{k }'s'}  \right] 
\end{multline*}
With vector triple product $    \A \times (\B \times \C) = \B(\A \cdot \C) - \C(\A \cdot \B)$,  the expression reads
\begin{multline*}
  \mathbf{E }\times \mathbf{B}
  = - \sum_{\mathbf{k },s} \sum_{\mathbf{k }',s'}  \frac{\hbar  }{2V} \sqrt{\frac{\mu_0 }{\epsilon_0 }\omega_k \omega_{k'}}
    \bigg[ \hat{a}_{\mathbf{k }s} \hat{a}_{\mathbf{k }'s'}  e ^{i(\phi_{\mathbf{k }} +\phi_{\mathbf{k }'})} 
  -\hat{a}_{\mathbf{k }s} \hat{a}^\dagger_{\mathbf{k }'s'} e ^{i(\phi_{\mathbf{k }}-\phi_{\mathbf{k }'})} \\ 
   -\hat{a}_{\mathbf{k }s} ^\dagger\hat{a}_{\mathbf{k }'s'} e ^{-i(\phi_{\mathbf{k }}-\phi_{\mathbf{k }'})}  
   +\hat{a}^\dagger_{\mathbf{k }s} \hat{a}^\dagger_{\mathbf{k }'s'} e ^{-i(\phi_{\mathbf{k }}+\phi_{\mathbf{k }'})} \bigg]
    \left[ \hat{\mathbf{k}}'\delta_{ss'}-\boldsymbol{\epsilon}_{\mathbf{k}'s'}\left( \boldsymbol{\epsilon}_{\mathbf{k }s}\cdot \hat{\mathbf{k}'} \right) \right] 
\end{multline*}
Under the spatial integration, the exponential with $\mathbf{k }+\mathbf{k}$ produce $\delta_{\mathbf{k },-\mathbf{k}}$ and $\mathbf{k }-\mathbf{k}$ produce $\delta_{\mathbf{k },\mathbf{k}}$.
As  such there are two patterns of the cross product.
\begin{equation*}
  \mathbf{E }\times \mathbf{B}  \propto
  \left[ \hat{a}_{\mathbf{k }s} \hat{a}_{-\mathbf{k }s}  \hat{a}^\dagger_{\mathbf{k }s}+ \hat{a}^\dagger_{-\mathbf{k }s} \right] \left[ -\mathbf{\hat{k}}' \right] 
\end{equation*}
This expression is odd, so summation over the same period cancel each other.
The other type, however 
\begin{equation*}
  \mathbf{E }\times \mathbf{B}  \propto
  \left[ -\hat{a}_{\mathbf{k }s} \hat{a}^\dagger_{\mathbf{k }s} - \hat{a}^\dagger_{\mathbf{k }s} \hat{a}_{\mathbf{k }s}  \right] \mathbf{\hat{k}}
\end{equation*}
do survive, since it is even.
We can also use the commutator relation to simplify further
\begin{equation*}
  \left[ -\hat{a}_{\mathbf{k }s} \hat{a}^\dagger_{\mathbf{k }s} - \hat{a}^\dagger_{\mathbf{k }s} \hat{a}_{\mathbf{k }s}    \hat{a}_{\mathbf{k }s}   \right] \hat{\mathbf{k}}
   = \left[ - 2\hat{a}^\dagger_{\mathbf{k }s \hat{a}_{\mathbf{k }s}  } -1\right] \mathbf{\hat{k}}
\end{equation*}
Also, the lasy term is odd, so it vanish under the summation.
Putting it all together
\begin{align*}
  \mathbf{P } = 
  \sum_{\mathbf{k },s } \hbar \omega_k \mu_0 \hat{a}^\dagger_{\mathbf{k }s} \hat{a}_{\mathbf{k }s}  \mathbf{\hat{k }}
  = \sum_{\mathbf{k },s}\hbar \mathbf{k } \hat{N }_{\mathbf{k }s} 
\end{align*}

\subsubsection{Linear and circular polarization}
Vector polarization $\boldsymbol{\epsilon}_{\mathbf{k }s}$ satisfy 
\begin{equation*}
  \boldsymbol{\epsilon}_{\mathbf{k }s} \cdot \boldsymbol{\epsilon}_{\mathbf{k }s'}=\delta_{ss'}
\end{equation*}
With linear vector polarization $\boldsymbol{\epsilon}_{\mathbf{k }s}$ and $\boldsymbol{\epsilon}_{\mathbf{k }s'}$, we maybe define circular polarization as 
\begin{equation*}
  \boldsymbol{\varepsilon}_{\mathbf{k }s} = \frac{1 }{\sqrt{2 }}\left( \boldsymbol{\epsilon}_{\mathbf{k }s}\pm i \boldsymbol{\epsilon}_{\mathbf{k }s'} \right) 
\end{equation*} 

\subsection{Second quantization of the Schrödinger equation.}
The time-dependent Schrödinger equation and time independent read as 
\begin{equation*}
i \frac{\partial \psi}{\partial t} = -\frac{1}{2m} \nabla^2 \psi + V(\mathbf{r})\psi \qquad
E_n \psi_n = -\frac{1}{2m} \nabla^2 \psi_n + V(\mathbf{r})\psi_n  
\end{equation*}
From separation of variables, we obtain solution called stationary state $\phi(\mathbf{r },t)=\phi_n (\mathbf{r})T(t)$.
Therefore, the general solution of the time-dependent Schrödinger equation is superposition of all separated solutions
\begin{equation*}
\psi(\mathbf{r}, t) = \sum_n a_n(t) \psi_n(\mathbf{r}) 
\end{equation*}
with $\dot{a}_n(t) = -i E_n a_n(t) $.

The energy $E$, or, more precisely, the expectation value of the Hamilton operator $\braket{H }$, can now be written as
\begin{equation*}
  \braket{H } 
  = \int d^3 r \, \psi^*(\mathbf{r}, t) \left[ -\frac{1}{2m} \nabla^2 + V(\mathbf{r}) \right] \psi(\mathbf{r}, t)
= \sum_{n} E_{n} a_{n}^{*} a_{n} = \sum_{n} \omega_{n} a_{n}^{*} a_{n}
\end{equation*}
Thus, it is consistent to interpret these result in terms of a second quantized Hamiltonian 
\begin{equation*}
  H = \sum_{n} E_{n} a_{n}^{\dagger} a_{n}
\end{equation*}
 
From the Heisenberg equation of motion
\begin{equation*}
  i \frac{dA}{dt  } = i [H,A]
\end{equation*}
we obtain equation of motion for the creation operators \(a_{n}\),
\begin{equation*}
  \frac{da_{n}}{dt} = i [H, a_{n}] = -i E_{n} a_{n}
\end{equation*}

Using the same argument, the second quantized Hamiltonian for fermion can be identified
\begin{equation*}
  H = \sum_{n} E_{n} c_{n}^{\dagger} c_{n}
\end{equation*}
and its Heisenberg equation of motion
\begin{equation*}
  \frac{dc_{n}}{dt} = i [H, c_{n}] = -i E_{n} c_{n}S
\end{equation*}
\end{document}